% Options for packages loaded elsewhere
\PassOptionsToPackage{unicode}{hyperref}
\PassOptionsToPackage{hyphens}{url}
\documentclass[
  english,
]{article}
\usepackage{xcolor}
\usepackage{amsmath,amssymb}
% section numbering enabled in preamble
\usepackage{iftex}
\ifPDFTeX
  \usepackage[T1]{fontenc}
  \usepackage[utf8]{inputenc}
  \usepackage{textcomp} % provide euro and other symbols
\else % if luatex or xetex
  \usepackage{unicode-math} % this also loads fontspec
  \defaultfontfeatures{Scale=MatchLowercase}
  \defaultfontfeatures[\rmfamily]{Ligatures=TeX,Scale=1}
\fi
\usepackage{lmodern}
\ifPDFTeX\else
  % xetex/luatex font selection
\fi
% Use upquote if available, for straight quotes in verbatim environments
\IfFileExists{upquote.sty}{\usepackage{upquote}}{}
\IfFileExists{microtype.sty}{% use microtype if available
  \usepackage[]{microtype}
  \UseMicrotypeSet[protrusion]{basicmath} % disable protrusion for tt fonts
}{}
\makeatletter
\@ifundefined{KOMAClassName}{% if non-KOMA class
  \IfFileExists{parskip.sty}{%
    \usepackage{parskip}
  }{% else
    \setlength{\parindent}{0pt}
    \setlength{\parskip}{6pt plus 2pt minus 1pt}}
}{% if KOMA class
  \KOMAoptions{parskip=half}}
\makeatother
\usepackage{longtable,booktabs,array}
\newcounter{none} % for unnumbered tables
\usepackage{calc} % for calculating minipage widths
% Correct order of tables after \paragraph or \subparagraph
\usepackage{etoolbox}
\makeatletter
\patchcmd\longtable{\par}{\if@noskipsec\mbox{}\fi\par}{}{}
\makeatother
% Allow footnotes in longtable head/foot
\IfFileExists{footnotehyper.sty}{\usepackage{footnotehyper}}{\usepackage{footnote}}
\makesavenoteenv{longtable}
\usepackage{graphicx}
\makeatletter
\newsavebox\pandoc@box
\newcommand*\pandocbounded[1]{% scales image to fit in text height/width
  \sbox\pandoc@box{#1}%
  \Gscale@div\@tempa{\textheight}{\dimexpr\ht\pandoc@box+\dp\pandoc@box\relax}%
  \Gscale@div\@tempb{\linewidth}{\wd\pandoc@box}%
  \ifdim\@tempb\p@<\@tempa\p@\let\@tempa\@tempb\fi% select the smaller of both
  \ifdim\@tempa\p@<\p@\scalebox{\@tempa}{\usebox\pandoc@box}%
  \else\usebox{\pandoc@box}%
  \fi%
}
% Set default figure placement to htbp
\def\fps@figure{htbp}
\makeatother
\ifLuaTeX
\usepackage[bidi=basic,shorthands=off]{babel}
\else
\usepackage[bidi=default,shorthands=off]{babel}
\fi
\ifLuaTeX
  \usepackage{selnolig} % disable illegal ligatures
\fi
\setlength{\emergencystretch}{3em} % prevent overfull lines
\providecommand{\tightlist}{%
  \setlength{\itemsep}{0pt}\setlength{\parskip}{0pt}}
% Appended by build_report_latex.py (after pandoc header). Tectonic-compatible.

\usepackage{xcolor}
\definecolor{ReportAccent}{HTML}{1F4E79}
\definecolor{ReportMuted}{HTML}{5B6770}
\definecolor{ReportLight}{HTML}{E8F1F8}

\PassOptionsToPackage{colorlinks=true,linkcolor=ReportAccent,urlcolor=ReportAccent,citecolor=ReportAccent}{hyperref}

\usepackage{float}
\usepackage[margin=2.2cm, headheight=15pt]{geometry}
\usepackage{caption}
\usepackage{subcaption}
\usepackage[most]{tcolorbox}
\usepackage{titlesec}
\usepackage{fancyhdr}
\usepackage{enumitem}

\floatplacement{figure}{H}

\graphicspath{{figures/}}

% Numbered sections with accent colour and rule
\titleformat{\section}
  {\normalfont\Large\bfseries\color{ReportAccent}}
  {\thesection}{0.75em}{}
  [\vspace{0.15em}{\color{ReportAccent!35}\titlerule[0.6pt]}]

\titleformat{\subsection}
  {\normalfont\large\bfseries\color{ReportAccent!85!black}}
  {\thesubsection}{0.6em}{}

\titlespacing*{\section}{0pt}{1.6em}{0.8em}
\titlespacing*{\subsection}{0pt}{1.2em}{0.5em}

% Figures
\captionsetup{
  font=small,
  labelfont={bf,color=ReportAccent},
  format=plain,
  labelsep=period,
  skip=8pt,
  singlelinecheck=false,
}

% Abstract panel
\newtcolorbox{abstractbox}{
  enhanced,
  colback=ReportLight,
  colframe=ReportAccent,
  boxrule=0.8pt,
  arc=2pt,
  left=12pt,
  right=12pt,
  top=10pt,
  bottom=10pt,
  before skip=0.5em,
  after skip=1.2em,
  before upper={\setlength{\parskip}{0.85em}},
}

% Headline statistics
\newtcolorbox{keystats}{
  enhanced,
  colback=ReportLight!45,
  colframe=ReportAccent!55,
  boxrule=0.5pt,
  arc=2pt,
  left=10pt,
  right=10pt,
  top=8pt,
  bottom=8pt,
  before skip=0.6em,
  after skip=1em,
}
\newenvironment{keystatslist}{%
  \begin{itemize}[leftmargin=1.2em, itemsep=0.25em, topsep=0.2em, parsep=0pt]
}{%
  \end{itemize}
}

% Compact scientific tables
\newenvironment{scitable}{%
  \par\vspace{0.4em}
  \begin{center}
  \footnotesize
  \setlength{\tabcolsep}{4pt}
  \renewcommand{\arraystretch}{1.25}
}{%
  \end{center}
  \par\vspace{0.6em}
}

\captionsetup[table]{font=small, labelfont=bf, skip=6pt}

% Running header / footer
\pagestyle{fancy}
\fancyhf{}
\fancyhead[L]{\small\color{ReportMuted}\nouppercase{\leftmark}}
\fancyhead[R]{\small\color{ReportMuted}\thepage}
\renewcommand{\headrulewidth}{0pt}
\renewcommand{\headrule}{\color{ReportAccent!30}\hrule width\headwidth height 0.4pt}
\fancypagestyle{plain}{%
  \fancyhf{}
  \fancyfoot[C]{\small\color{ReportMuted}\thepage}
  \renewcommand{\headrulewidth}{0pt}
}

\newcommand{\um}{\,\text{$\mu$m}}
\newcommand{\umsq}{\,\text{$\mu$m}^2}

\setcounter{secnumdepth}{2}
\usepackage{bookmark}
\IfFileExists{xurl.sty}{\usepackage{xurl}}{} % add URL line breaks if available
\urlstyle{same}
\hypersetup{
  pdftitle={Cell Classifier and Cell Size Report},
  pdflang={en},
  hidelinks,
  pdfcreator={LaTeX via pandoc}}

\title{Cell Classifier and Cell Size Report}
\author{}
\date{23 May 2026}

\begin{document}
\begin{titlepage}
\thispagestyle{empty}
\centering
\vspace*{2.2cm}
{\color{ReportAccent}\rule{\textwidth}{2pt}\par}
\vspace{1.4cm}
{\Huge\bfseries\color{ReportAccent} Cell Classifier and Cell Size Report\par}
\vspace{1.2cm}
{\large\color{ReportMuted} Automatic quality filtering and morphometry across frogs\par}
\vfill
{\large 23 May 2026\par}
\vspace{2cm}
{\color{ReportAccent}\rule{\textwidth}{2pt}\par}
\end{titlepage}
\setcounter{page}{1}

\section{Summary}\label{summary}
\begin{abstractbox}
We trained a deep-learning model (ResNet18 backbone) that looks at every
segmented cell and decides whether it is a \textbf{good} cell worth
measuring or a \textbf{bad} segmentation that should be thrown away. On
an independent test set of \textbf{1,300} cells the model performs
strongly: it correctly accepts good cells \textbf{94.71\%} of the time
and correctly discards bad cells \textbf{98.65\%} of the time, with
overall accuracy \textbf{98.08\%}. With \textbf{selective rejection}
(allowing the model to defer uncertain cells), the accepted subset
reaches \textbf{98.67\%} accuracy at \textbf{98.00\%} coverage.

Applying this filter to the full imaging dataset yields \textbf{57,520}
good cells from \textbf{479} frogs for morphometry. The typical (median)
cell covers \textbf{290.69 \(\mu\)m\(^2\)}, with median N/C ratio
\textbf{0.13}. Per-frog mean cell area spans \textbf{212.55--482.23
\(\mu\)m\(^2\)}. Nucleus--cell scaling is negatively allometric (nucleus
grows slower than cell area; larger cells are relatively less
nucleus-dense).
\end{abstractbox}

\section{Part A --- Baseline classifier (single decision
threshold)}\label{part-a-baseline-classifier-single-decision-threshold}

\subsubsection{How it decides}\label{how-it-decides}

The classifier outputs a probability \texttt{p(good)} between 0 and 1
for every cell. In \textbf{baseline mode} we keep the cell if
\textbf{p(good) \(\geq\) 0.50} and discard it otherwise. No cells are
deferred.

\subsubsection{Headline metrics on the test
split}\label{headline-metrics-on-the-test-split}

{\def\LTcaptype{none} % do not increment counter
\begin{scitable}
\begin{longtable}[]{@{}ll@{}}
\toprule\noalign{}
Metric & Value \\
\midrule\noalign{}
\endhead
\bottomrule\noalign{}
\endlastfoot
Accuracy & \textbf{98.08\%} \\
Recall (sensitivity to good cells) & \textbf{94.71\%} \\
Precision & \textbf{92.27\%} \\
F1 & \textbf{93.47\%} \\
Specificity (correctly discarded bad cells) & \textbf{98.65\%} \\
\end{longtable}
\end{scitable}
}

\subsubsection{Confusion matrix --- raw
counts}\label{confusion-matrix-raw-counts}

{\def\LTcaptype{none} % do not increment counter
\begin{scitable}
\begin{longtable}[]{@{}llll@{}}
\toprule\noalign{}
True ~Pred & Pred bad & Pred good & Row total \\
\midrule\noalign{}
\endhead
\bottomrule\noalign{}
\endlastfoot
True bad & \textbf{1096} (TN) & \textbf{15} (FP) & 1111 \\
True good & \textbf{10} (FN) & \textbf{179} (TP) & 189 \\
\end{longtable}
\end{scitable}
}

\subsubsection{Confusion matrix --- row-normalized
(\%)}\label{confusion-matrix-row-normalized}

{\def\LTcaptype{none} % do not increment counter
\begin{scitable}
\begin{longtable}[]{@{}lll@{}}
\toprule\noalign{}
True ~Pred & Pred bad & Pred good \\
\midrule\noalign{}
\endhead
\bottomrule\noalign{}
\endlastfoot
True bad & 98.65\% & 1.35\% \\
True good & 5.29\% & 94.71\% \\
\end{longtable}
\end{scitable}
}

\textbf{Interpretation.} When the model says a cell is good, it is right
about \textbf{92.27\%} of the time. About \textbf{94.71\%} of truly good
cells are kept.

\subsubsection{Qualitative examples --- what the model is
seeing}\label{qualitative-examples-what-the-model-is-seeing}

Each panel shows up to 20 randomly sampled test cells from that
category, with the model's \texttt{p(good)} printed above each cell.

\textbf{Correctly kept good cells (true positives):}

\begin{figure}[H]
\centering
\includegraphics[width=0.98\linewidth]{./figures/TruePositive.png}
\caption{True Positive examples}
\end{figure}

\textbf{Correctly discarded bad cells (true negatives):}

\begin{figure}[H]
\centering
\includegraphics[width=0.98\linewidth]{./figures/TrueNegative.png}
\caption{True Negative examples}
\end{figure}

\textbf{Bad cells kept by mistake (false positives):}

\begin{figure}[H]
\centering
\includegraphics[width=0.98\linewidth]{./figures/FalsePositive.png}
\caption{False Positive examples}
\end{figure}

\textbf{Good cells rejected by mistake (false negatives):}

\begin{figure}[H]
\centering
\includegraphics[width=0.98\linewidth]{./figures/FalseNegative.png}
\caption{False Negative examples}
\end{figure}

\section{Part B --- Adding a ``not sure'' option (selective
rejection)}\label{part-b-adding-a-not-sure-option-selective-rejection}

Baseline mode always produces an answer. For biology it is sometimes
safer to defer the most ambiguous crops.

\subsubsection{How it decides}\label{how-it-decides-1}

We split the probability axis into three zones using two thresholds
(after temperature scaling with T = \textbf{1.8838}):

\begin{keystats}
\begin{itemize}
\tightlist
\item
  \textbf{p(good) \(\leq\) 0.1} \(\rightarrow\) \textbf{bad} (accepted
  as bad)
\item
  \textbf{p(good) \(\geq\) 0.76} \(\rightarrow\) \textbf{good} (accepted
  as good)
\item
  \textbf{0.1 \textless{} p(good) \textless{} 0.76} \(\rightarrow\)
  \textbf{rejected / unsure} (deferred)
\end{itemize}
\end{keystats}

\subsubsection{Headline metrics --- baseline vs selective (test
split)}\label{headline-metrics-baseline-vs-selective-test-split}

All selective-rejection metrics below are computed \textbf{only on cells
the model accepted.}

{\def\LTcaptype{none} % do not increment counter
\begin{scitable}
\begin{longtable}[]{@{}lll@{}}
\toprule\noalign{}
Metric & Baseline & Selective rejection (accepted only) \\
\midrule\noalign{}
\endhead
\bottomrule\noalign{}
\endlastfoot
Coverage (cells that got a decision) & 100.00\% & \textbf{98.00\%} \\
Accuracy & 98.08\% & \textbf{98.67\%} \\
Recall & 94.71\% & \textbf{97.27\%} \\
Precision & 92.27\% & \textbf{93.68\%} \\
F1 & 93.47\% & \textbf{95.44\%} \\
Specificity & 98.65\% & \textbf{98.90\%} \\
False positives (kept by mistake) & 15 & \textbf{12} \\
False negatives (missed good cells) & 10 & \textbf{5} \\
\end{longtable}
\end{scitable}
}

\subsubsection{Coverage and rejection
summary}\label{coverage-and-rejection-summary}

\begin{itemize}
\tightlist
\item
  Accepted: \textbf{1,274 / 1,300} cells (\textbf{98.00\%})
\item
  Rejected / unsure: \textbf{26 / 1,300}

  \begin{itemize}
  \tightlist
  \item
    Among rejected: \textbf{6} truly good, \textbf{20} truly bad
  \end{itemize}
\end{itemize}

\textbf{Intuition.} False negatives drop from \textbf{10} to \textbf{5}
under selective rejection; false positives from \textbf{15} to
\textbf{12}.

\subsubsection{Accepted-only confusion matrix ---
counts}\label{accepted-only-confusion-matrix-counts}

{\def\LTcaptype{none} % do not increment counter
\begin{scitable}
\begin{longtable}[]{@{}llll@{}}
\toprule\noalign{}
True ~Pred (accepted) & Pred bad & Pred good & Row total \\
\midrule\noalign{}
\endhead
\bottomrule\noalign{}
\endlastfoot
True bad & \textbf{1079} (TN) & \textbf{12} (FP) & 1091 \\
True good & \textbf{5} (FN) & \textbf{178} (TP) & 183 \\
\end{longtable}
\end{scitable}
}

\subsubsection{Accepted-only confusion matrix --- row-normalized
(\%)}\label{accepted-only-confusion-matrix-row-normalized}

{\def\LTcaptype{none} % do not increment counter
\begin{scitable}
\begin{longtable}[]{@{}lll@{}}
\toprule\noalign{}
True ~Pred (accepted) & Pred bad & Pred good \\
\midrule\noalign{}
\endhead
\bottomrule\noalign{}
\endlastfoot
True bad & 98.90\% & 1.10\% \\
True good & 2.73\% & 97.27\% \\
\end{longtable}
\end{scitable}
}

\subsubsection{Qualitative examples under the selective
policy}\label{qualitative-examples-under-the-selective-policy}

\begin{figure}[H]
\centering
\includegraphics[width=0.98\linewidth]{./figures/AcceptedGood.png}
\caption{Accepted good examples}
\end{figure}

\begin{figure}[H]
\centering
\includegraphics[width=0.98\linewidth]{./figures/AcceptedBad.png}
\caption{Accepted bad examples}
\end{figure}

\begin{figure}[H]
\centering
\includegraphics[width=0.98\linewidth]{./figures/RejectedUncertain.png}
\caption{Rejected uncertain examples}
\end{figure}

\begin{figure}[H]
\centering
\includegraphics[width=0.98\linewidth]{./figures/AcceptedFalsePositive.png}
\caption{Accepted false-positive examples}
\end{figure}

\begin{figure}[H]
\centering
\includegraphics[width=0.98\linewidth]{./figures/AcceptedFalseNegative.png}
\caption{Accepted false-negative examples}
\end{figure}

\begin{figure}[H]
\centering
\includegraphics[width=0.98\linewidth]{./figures/Threshold.png}
\caption{Threshold histogram}
\end{figure}

\subsubsection{What improved}\label{what-improved}

\begin{keystats}
\begin{itemize}
\tightlist
\item
  \textbf{Recall} improved (\textbf{94.71\%} \(\rightarrow\)
  \textbf{97.27\%}) on accepted cells.
\item
  \textbf{Precision} improved (\textbf{92.27\%} \(\rightarrow\)
  \textbf{93.68\%}).
\item
  \textbf{Accuracy} reached \textbf{98.67\%} on the accepted subset at
  \textbf{98.00\%} coverage.
\end{itemize}
\end{keystats}

\section{Part C --- Biological results on the full
dataset}\label{part-c-biological-results-on-the-full-dataset}

After running the model with selective rejection on the full imaging
dataset:

\begin{keystats}
\begin{itemize}
\tightlist
\item
  \textbf{383,915} detections were scored across all images.
\item
  \textbf{57,520} were accepted as \textbf{good} and measured
  (\textbf{15.0\%} yield).
\item
  \textbf{57,218} good cells have a matched nucleus for N/C and
  regression analyses.
\end{itemize}
\end{keystats}

\subsection{C1. Cell-area distribution (good cells
only)}\label{c1.-cell-area-distribution-good-cells-only}

\begin{figure}[H]
\centering
\includegraphics[width=0.98\linewidth]{./figures/AreaDistribution.png}
\caption{Cell area distribution}
\end{figure}

\begin{keystats}
\begin{itemize}
\tightlist
\item
  \textbf{Median cell area: 290.69 \(\mu\)m\(^2\)} (IQR 257.70--335.48
  \(\mu\)m\(^2\)).
\item
  \textbf{Mean \(\pm\) SD: 303.10 \(\pm\) 67.07 \(\mu\)m\(^2\).}
\end{itemize}
\end{keystats}

\subsection{C2. Cell-diameter distribution (good cells
only)}\label{c2.-cell-diameter-distribution-good-cells-only}

\begin{figure}[H]
\centering
\includegraphics[width=0.98\linewidth]{./figures/DiameterDistribution.png}
\caption{Cell diameter distribution}
\end{figure}

\begin{keystats}
\begin{itemize}
\tightlist
\item
  \textbf{Long axis (major): 24.52 \(\mu\)m} (IQR 22.79--26.78
  \(\mu\)m).
\item
  \textbf{Short axis (minor): 15.23 \(\mu\)m} (IQR 14.16--16.44
  \(\mu\)m).
\item
  \textbf{Axis ratio (major \(\div\) minor): median 1.62}.
\end{itemize}
\end{keystats}

\subsection{C3. Nucleus size
distribution}\label{c3.-nucleus-size-distribution}

\begin{figure}[H]
\centering
\includegraphics[width=0.98\linewidth]{./figures/NucleusDistribution.png}
\caption{Nucleus distribution}
\end{figure}

\begin{keystats}
\begin{itemize}
\tightlist
\item
  \textbf{Nucleus area: 39.06 \(\mu\)m\(^2\)} (IQR 34.79--46.13
  \(\mu\)m\(^2\)).
\item
  \textbf{Long / short axis medians: 8.96 / 5.66 \(\mu\)m}.
\end{itemize}
\end{keystats}

\subsection{C4. Nucleus-to-cell area ratio (N/C
ratio)}\label{c4.-nucleus-to-cell-area-ratio-nc-ratio}

\begin{figure}[H]
\centering
\includegraphics[width=0.98\linewidth]{./figures/NCRatioDistribution.png}
\caption{N/C ratio distribution}
\end{figure}

\begin{keystats}
\begin{itemize}
\tightlist
\item
  \textbf{Median N/C ratio: 0.13} (about 13\% of cell area).
\item
  \textbf{Mean \(\pm\) SD: 0.14 \(\pm\) 0.04}.
\end{itemize}
\end{keystats}

\subsection{C5. Per-frog variation in mean cell
area}\label{c5.-per-frog-variation-in-mean-cell-area}

\begin{figure}[H]
\centering
\includegraphics[width=0.98\linewidth]{./figures/PerFrogBoxplot.png}
\caption{Per-frog cell-area boxplot}
\end{figure}

Per-frog boxplot of cell area (top-40 frogs by cell count, sorted by
per-frog median cell area). Per-frog mean cell area spans
\textbf{212.55--482.23 \(\mu\)m\(^2\)}; typical frog's average
\textbf{290.34 \(\mu\)m\(^2\)}.

\begin{figure}[H]
\centering
\includegraphics[width=0.98\linewidth]{./figures/PerFrogNucleusBoxplot.png}
\caption{Per-frog nucleus-area boxplot}
\end{figure}

Same frogs and ordering for nucleus area.

\subsection{C6. Does the nucleus scale with cell size?
(regression)}\label{c6.-does-the-nucleus-scale-with-cell-size-regression}

\begin{figure}[H]
\centering
\includegraphics[width=0.98\linewidth]{./figures/NucleusVsCellRegression.png}
\caption{Nucleus area vs cell area regression}
\end{figure}

We fit a regression in \textbf{log--log space}. A slope of \textbf{1} is
isometric scaling; \textbf{below 1} is negative allometry.
\textbf{Pooled OLS slope: 0.526} (R\(^2\) = 0.255);
\textbf{mixed-effects slope: 0.149} (random intercept per frog). Scaling
is negatively allometric (nucleus grows slower than cell area; larger
cells are relatively less nucleus-dense).

\section{Discussion}\label{discussion}

This report combines \textbf{classifier validation} on a held-out test
set with \textbf{morphometry} on the full cohort after selective
rejection. On the test split, baseline accuracy is \textbf{98.08\%};
selective rejection improves accepted-only accuracy to \textbf{98.67\%}
while deferring \textbf{26} ambiguous cells (\textbf{98.00\%} coverage).
That trade-off is appropriate when morphometry should describe only
cells the model judged with confidence.

On the full dataset, \textbf{15.0\%} of scored detections were accepted
as good (median per-frog yield \textbf{17.7\%}). Morphometric summaries
therefore describe the \textbf{accepted subset}, not every segmented
object. Frogs with low yield may be under-represented.

Cell size varies substantially between frogs (ICC \textbf{0.661};
per-frog mean area \textbf{212.55--482.23 \(\mu\)m\(^2\)}).
Nucleus--cell scaling is negatively allometric (nucleus grows slower
than cell area; larger cells are relatively less nucleus-dense): larger
cells tend to have proportionally smaller nuclei (lower N/C). The mixed
slope (\textbf{0.149}) is smaller than the pooled OLS slope
(\textbf{0.526}), reflecting that much nucleus--cell correlation is
\textbf{between frogs}, not only within them.

\section{Glossary}\label{glossary}

{\def\LTcaptype{none} % do not increment counter
\begin{scitable}
\begin{longtable}[]{@{}
  >{\raggedright\arraybackslash}p{(\linewidth - 2\tabcolsep) * \real{0.4000}}
  >{\raggedright\arraybackslash}p{(\linewidth - 2\tabcolsep) * \real{0.6000}}@{}}
\toprule\noalign{}
\begin{minipage}[b]{\linewidth}\raggedright
Term
\end{minipage} & \begin{minipage}[b]{\linewidth}\raggedright
Meaning
\end{minipage} \\
\midrule\noalign{}
\endhead
\bottomrule\noalign{}
\endlastfoot
\textbf{Good cell} & A segmented cell accepted by the quality classifier
for morphometry. \\
\textbf{Bad cell} & A segmentation rejected as poor quality. \\
\textbf{Selective rejection} & Three-way policy: accept as good, accept
as bad, or defer as uncertain. \\
\textbf{Coverage} & Fraction of test cells that received a good/bad
decision (not deferred). \\
\textbf{Accuracy} & Fraction of all labelled cells classified
correctly. \\
\textbf{Recall} & Of truly good cells, fraction the model kept. \\
\textbf{Precision} & Of cells called good, fraction that are truly
good. \\
\textbf{Specificity} & Of truly bad cells, fraction the model
discarded. \\
\textbf{F1} & Harmonic mean of precision and recall. \\
\textbf{N/C ratio} & Nucleus area divided by cell area (unitless). \\
\textbf{Yield / good fraction} & Share of scored detections accepted as
good. \\
\textbf{ICC} & Share of cell-area variation explained by which frog a
cell came from. \\
\textbf{Allometry} & How nucleus size scales with cell size; slope = 1
is isometric, \textless{} 1 is negative allometry. \\
\textbf{Temperature scaling} & Post-hoc calibration of classifier
probabilities. \\
\end{longtable}
\end{scitable}
}

\end{document}
